% stylelibrary v2.3.1 (transparent box backgrounds by default: \opacityBack=0; see README.md).
\usepackage[many]{tcolorbox} %box environment
\usepackage[normalem]{ulem} %custom underline

% Bar width and spacing
\newlength{\framelinewidth}
\setlength{\framelinewidth}{0.6mm}
\newlength{\marginsep}
\setlength{\marginsep}{3pt}
\newlength{\defparindent}
\setlength{\defparindent}{15pt}
\setlength{\parindent}{\defparindent}
\newlength{\boxparskip}
\setlength{\boxparskip}{\parskip}

%Define colors
\definecolor{aicolor}{HTML}{038DFC}
\definecolor{editedaicolor}{HTML}{2B65C6}
\definecolor{humancolor}{HTML}{22AA22}
\definecolor{humanaicolor}{HTML}{5700FF}
\definecolor{defaultcolback}{HTML}{FFFFFF}  % white
\newcommand{\opacityFrame}{0.5} %Opacity Percentage for the Frame
\newcommand{\opacityText}{0.7} %Opacity Percentage for the Label Text
\newcommand{\opacityBack}{0} %Opacity Percentage for the Background (0 = fully transparent)

%Define footenote commands: Creates an aux file for the footnotes inside the boxes
\newif\iffootaux
\newcounter{footaux}% for \csname
\newcounter{auxpage}% to handle frontmatter etc.
\setcounter{auxpage}{1}

\makeatletter

\def\tcb@footnote@setup{
\let\footnotetext=\footaux
\let\footnote=\footnoteaux
}

\def\footnoteaux{\@ifnextchar[\@xfootnoteaux{\stepcounter{footnote}%
  \footnotemark[\thefootnote]\@footaux[\thefootnote]}}

\def\@xfootnoteaux[#1]{\footnotemark[#1]\@footaux[#1]}

\def\footaux{\@ifnextchar[\@footaux{\@footaux[\thefootnote]}}
\def\@footaux[#1]#2{\protected@write{\@auxout}{\let\theauxpage\relax}{\string\newfootnote{\theauxpage}{#1}{#2}}}

\newcommand{\newfootnote}[3]{% #1 = \theauxpage, #2 = \thefoontnote, #3 = text
  \ifnum#1=1\relax
    \footnotetext[#2]{#3}%
  \else
    \stepcounter{footaux}%
    \expandafter\xdef\csname footauxpage\thefootaux\endcsname{#1}%
    \expandafter\xdef\csname footauxnumber\thefootaux\endcsname{#2}%
    \expandafter\protected@xdef\csname footauxtext\thefootaux\endcsname{#3}%
  \fi}

% to run at the end of every page (debug version)
\newcommand{\addfootaux}{\loop
  \@ifundefined{footauxpage\thefootaux}{\footauxfalse}{%
    \count1=\csname footauxpage\thefootaux\endcsname\relax
    \advance\count1 by -1
    \ifnum\count1=\value{auxpage}%
      \expandafter\let\expandafter\footauxnumber\csname footauxnumber\thefootaux\endcsname
      \expandafter\let\expandafter\footauxtext\csname footauxtext\thefootaux\endcsname
      \footnotetext[\footauxnumber]{\footauxtext}%
      \footauxtrue
    \else
      \footauxfalse
    \fi}
  \iffootaux
    \global\expandafter\let\csname footauxpage\thefootaux\endcsname\relax
    \global\expandafter\let\csname footauxnumber\thefootaux\endcsname\relax
    \global\expandafter\let\csname footauxtext\thefootaux\endcsname\relax
    \stepcounter{footaux}%
  \repeat
\stepcounter{auxpage}}
\makeatother

\AtBeginDocument{\setcounter{footaux}{1}}% reset counter after \newfootnote
\AddToHook{shipout/after}{\addfootaux}


%Define box styles
\makeatletter
\tcbset{
    roundbox/.style={
    % Box Options
    rounded corners,
    arc=3mm,
    boxrule=\framelinewidth
  },
  rectanglebox/.style={
    % Box Options
    sharp corners,
    boxrule=\framelinewidth
  },
    linemargin/.style={
    % Box Options
    % Make frame with 1 line
    sharp corners,
    leftrule=\framelinewidth,    % Thickness of the left line
    rightrule=-0.1mm,            % Removes right frame
    toprule=-0.1mm,              % Removes top frame
    bottomrule=-0.125mm,         % Removes bottom frame
    % % make white background
    colback=white,
    opacityback=\opacityBack,
    % Oversize makes bigger, so that text alings with rest of the document
    oversize,
    % box margins
    left=\marginsep,
    right=\marginsep,
    top=0.5pt,
    bottom=0.5pt,
    before skip=0.5\baselineskip,
    after skip=0.5\baselineskip
  },
  alignoptions/.style={
    % box margins
    left=\marginsep,
    right=\marginsep,
    top=0.5pt,
    bottom=0.5pt,
    before skip=0.5\baselineskip,
    after skip=\parskip,
    before upper={
    \parindent=\defparindent
    \setlength{\parskip}{\boxparskip}
    }
  },
  %Format for footnotes to appear in the correct page
  footnoteformat/.style={
  before upper app=\tcb@footnote@setup
  }
}
\makeatother

\tcbset{
  enhancebreak/.style={
    % Box Options
    enhanced,
    breakable,
    skin first=enhanced,
    skin middle=enhanced,
    skin last=enhanced,
    % Oversize makes bigger, so that text alings with rest of the document
    oversize,
  },
}

% Label-to-bar separation (default 0.7pt; see README for usage).
% The overlay adds the fixed 2.65pt node-edge-to-glyph inset
% (default inner sep at \tiny) to this value for the node xshift.
\newcommand{\labelsepvalue}{0.7pt}
% `labeltopleft` (off) puts the label horizontally at the top-left
% corner, at the start of the block, instead of the rotated label.
\newif\iflabeltopleft
\tcbset{labeltopleft/.is if=labeltopleft}
% Per-box label text and color (set by each box env).
\newcommand{\boxlabeltext}{}
\newcommand{\boxlabelcolor}{black}
\tcbset{
  boxlabel/.store in=\boxlabeltext,
  boxlabelcolor/.store in=\boxlabelcolor,
  labeltopleftlayout/.code={%
    \iflabeltopleft \tcbset{top=12pt} \fi
  }
}
\tcbset{labelsep/.store in=\labelsepvalue}

% `extendbar` (on) extends the bar past the label glyphs by
% `barextend` (default 1pt); see README for usage. The floor is the
% label-node width + barextend − 7.19pt (calibrated node-to-glyph offset).
\newif\ifextendbar
\extendbartrue
\tcbset{extendbar/.is if=extendbar}
\newcommand{\barextendvalue}{1pt}
\tcbset{barextend/.store in=\barextendvalue}
\newcommand{\barextendoffset}{7.19pt}
\newsavebox{\extendbarbox}
\tcbset{
  extendbarmeasure/.code={%
    \iflabeltopleft\else
    \ifextendbar
      \sbox{\extendbarbox}{\tikz{\node[font=\tiny\bfseries\sffamily, align=right]{#1};}}%
      \tcbset{height=\dimexpr\wd\extendbarbox + \barextendvalue - \barextendoffset\relax, height plus=10000pt}%
    \fi
    \fi
  }
}

% Define Tcolorbox environments
%% Args: #1 additional options, #2 color, #3 Side text label
\newtcolorbox{labeledblock}[3][]{
    colframe=#2,opacityframe=\opacityFrame,
    enhancebreak, alignoptions, linemargin,
    colback=white,
    frame hidden,
    borderline west={\framelinewidth}{0pt}{#2, opacity=\opacityFrame},
    overlay={
    \iflabeltopleft
    \node[anchor=north west, inner sep=0pt, font=\scriptsize\bfseries\sffamily, text=\boxlabelcolor, text opacity=\opacityText] at ([xshift=7.5pt, yshift=-2pt]frame.north west) {\boxlabeltext};
    \else
    \node[
    #2, font=\tiny\bfseries\sffamily,
    rotate=90, anchor=south east, align=right,
    text opacity=\opacityText
      ]
      at ([xshift=2.65pt-\labelsepvalue, yshift=3.8pt]frame.north west) {#3};
    \fi
    }, footnoteformat,
    boxlabel={#3}, boxlabelcolor=#2,
    #1,
    labeltopleftlayout, extendbarmeasure={#3}
}

\newtcolorbox{humanblock}[1][]{
    colframe=humancolor, opacityframe=\opacityFrame,
    enhancebreak, alignoptions, linemargin,
    colback=white,
    frame hidden,
    borderline west={\framelinewidth}{0pt}{humancolor, opacity=\opacityFrame},
    overlay={
    \iflabeltopleft
    \node[anchor=north west, inner sep=0pt, font=\scriptsize\bfseries\sffamily, text=\boxlabelcolor, text opacity=\opacityText] at ([xshift=7.5pt, yshift=-2pt]frame.north west) {\boxlabeltext};
    \else
    \node[
    humancolor, font=\tiny\bfseries\sffamily,
    rotate=90, anchor=south east, align=right,
    text opacity=\opacityText
      ]
      at ([xshift=2.65pt-\labelsepvalue, yshift=3.8pt]frame.north west) {Human};
    \fi
    }, footnoteformat,
    boxlabel=Human, boxlabelcolor=humancolor,
    #1,
    labeltopleftlayout, extendbarmeasure={Human}
}

\newtcolorbox{humanaiblock}[1][]{
    colframe=humanaicolor, opacityframe=\opacityFrame,
    enhancebreak, alignoptions, linemargin,
    colback=white,
    frame hidden,
    borderline west={\framelinewidth}{0pt}{humanaicolor, opacity=\opacityFrame},
    overlay={
    \iflabeltopleft
    \node[anchor=north west, inner sep=0pt, font=\scriptsize\bfseries\sffamily, text=\boxlabelcolor, text opacity=\opacityText] at ([xshift=7.5pt, yshift=-2pt]frame.north west) {\boxlabeltext};
    \else
    \node[humanaicolor, font=\tiny\bfseries\sffamily, rotate=90, anchor=south east, align=right,
      text opacity=\opacityText
      ]
      at ([
      xshift=2.65pt-\labelsepvalue, yshift=3.8pt
      ]frame.north west) {Human + AI};
    \fi
    }, footnoteformat,
    boxlabel={Human + AI}, boxlabelcolor=humanaicolor,
    #1,
    labeltopleftlayout, extendbarmeasure={Human + AI}
}

\newtcolorbox{aiblock}[1][]{
    colframe=aicolor, opacityframe=\opacityFrame,
    enhancebreak, alignoptions, linemargin,
    colback=white,
    frame hidden,
    borderline west={\framelinewidth}{0pt}{aicolor, opacity=\opacityFrame},
    overlay={
    \iflabeltopleft
    \node[anchor=north west, inner sep=0pt, font=\scriptsize\bfseries\sffamily, text=\boxlabelcolor, text opacity=\opacityText] at ([xshift=7.5pt, yshift=-2pt]frame.north west) {\boxlabeltext};
    \else
    \node[aicolor, font=\tiny\bfseries\sffamily, rotate=90, anchor=south east, align=right,
      text opacity=\opacityText
      ]
      at ([xshift=2.65pt-\labelsepvalue, yshift=3.8pt]frame.north west) {AI};
    \fi
    }, footnoteformat,
    boxlabel=AI, boxlabelcolor=aicolor,
    #1,
    labeltopleftlayout, extendbarmeasure={AI}
}

\newtcolorbox{editedaiblock}[1][]{
    colframe=editedaicolor, opacityframe=\opacityFrame,
    enhancebreak, alignoptions, linemargin,
    colback=white,
    frame hidden,
    borderline west={\framelinewidth}{0pt}{editedaicolor, opacity=\opacityFrame},
    overlay={
    \iflabeltopleft
    \node[anchor=north west, inner sep=0pt, font=\scriptsize\bfseries\sffamily, text=\boxlabelcolor, text opacity=\opacityText] at ([xshift=7.5pt, yshift=-2pt]frame.north west) {\boxlabeltext};
    \else
    \node[editedaicolor, font=\tiny\bfseries\sffamily, rotate=90, anchor=south east, align=right,
      text opacity=\opacityText
      ]
      at ([
      xshift=2.65pt-\labelsepvalue, yshift=3.8pt]frame.north west) {Edited AI};
    \fi
    }, footnoteformat,
    boxlabel={Edited AI}, boxlabelcolor=editedaicolor,
    #1,
    labeltopleftlayout, extendbarmeasure={Edited AI}
}


%%%---------------------- Underline Commands -------------------------

% Dash pattern:
\gdef\dashpattern{\kern.05em\lower.3ex\vbox{\hrule width.2em\kern.1ex}}

\def\humanul{\textcolor{humancolor}\bgroup\markoverwith{\kern.05em\lower.3ex\vbox{\hrule width.2em\kern.1ex}\kern.05em}\ULon}

\def\aiul{\textcolor{aicolor}\bgroup\markoverwith{\kern.05em\lower.3ex\vbox{\hrule width.2em\kern.1ex}\kern.05em}\ULon}
